\chapter{Computation of Euler-Lagrange solutions}\label{sec:discretization-euler-lagrange}

Perhaps the most straightforward approach to compute numerically solutions that extremize the action given a Lagrangian is to compute the ODE that results from Euler-Lagrange (\autoref{eq:euler-lagrange}), and then use some numerical ODE computation scheme based on the underlying type of manifold. For example if $Q = \RR^n$ (or is, in general, a vector space) then we can use Runge-Kutta, or if $Q$ is a Lie group we could use Runge-Kutta-Kaan.

However, discretization the resulting ODE based on $q(t)$ can have undesirable effects due to the loss of the geometric information of the configuration manifold $Q$. These will be explained in \autoref{sec:pitfalls-discrete-el-ode}.

Instead, there is an alternative approach of discretizing the tangent bundle of the configuration space $TQ$ as $Q \times Q$, and deriving discrete analogs of the Euler-Lagrange equations themselves. This approach, called \newterm{variational integration} \todo{is it?} will be explained in \autoref{sec:discretizing-configuration-space}

But, to understand the tradeoffs between the approaches, we first need to examine some important geometric properties Euler-Lagrange.

\todo{Think about how much I want to explain these things\ldots}

\section{Geometric properties Euler-Lagrange}

These are discussed more in depth in~\cite[Chapter 1.2]{marsdenDiscreteMechanicsVariational2001}.

\subsubsection{Simpleticity}

Given a Lagrangian $L: TQ \to \RR$, we can define the \newterm{Lagrangian vector field} $X_L: TQ \to T(TQ)$ as the second order vector field that satisfies:

\[
	D_{EL} L \circ X_L = 0
\]

\subsubsection{Preservation of momentum maps}


\section{Pitfalls of discretizing the Euler-Lagrange ODE}\label{sec:pitfalls-discrete-el-ode}

The fundamental problem of integrating numerically the Euler-Lagrange ODE is, essentially, that we lose the geometry of the configuration manifold $Q$. The numerical inaccuracies can ``go out'' of what should be possible in $Q$ and keep degenerating since there is no correction mechanism.

Let's look at an example.

\todo{Talk about how energy goes down.}

\section{Discretizing the configuration space}\label{sec:discretizing-configuration-space}

The discrete analog of continuous paths $q(t) : \RR \to TQ$ are sequences of points in $Q$. Generally, we like to work  Then, the discrete analog of the Lagrangian $L_d$ takes two consecutive points instead of one point in $TQ$.

\[ L: TQ \to \RR \Rightarrow L_d : Q \times Q \to \RR \]

Then, the discrete action is now composed of sums of these discrete Lagrangians. That is, for a path $\mathbb{q} = {\{ q_i \in Q \}}_{i=0}^{N-1}$, the discrete action $S_d$ is:

\begin{equation}
	S_d (\mathbf{q}) = \sum_{k=0}^{N-2} \Ld(q_k, q_{k+1})
	\label{eq:action-discrete}
\end{equation}

If we define an \newterm{exact discrete Lagrangian} $\Ld^e$ as
\begin{equation}
	L_d^e (q_0, q_1) = \int_0^h L(q_{0, 1}(t), \dot{q}_{0, 1}(t)) \odif{t}
\end{equation}
where $q_{0, 1}$ is the path that minimizes the action along all the paths that connect $q_0$ and $q_1$, then the discrete formulation gives identical results to the discrete one.

Of course, finding the path $q_{0,1}$ is as hard as the original continuous problem; but in practice we often define the discrete Lagrangian $\Ld$ using some quadrature rule that approximates the integral of the exact discrete Lagrangian.

If we apply Hamilton's principle to the discrete action (that is, $\nabla S_d = 0$ \todo{Maybe I should write out the derivation}), we find the discrete Euler-Lagrange (DEL) equations:
\begin{equation}
	\mdif{2} L_d (q_{k-1}, q_k) + \mdif{1} L_d (q_k, q_{k+1}) = 0
	\label{eq:euler-lagrange-discrete}\
	\tag{DEL}
\end{equation}

These equations are well-defined for any manifold, since both $\mdif{2} L_d (q_{k-1}, q_k)$ and $\mdif{1} L_d (q_k, q_{k+1})$ live in $T_{q_k} Q$, and the local tangent space of any smooth manifold always forms a vector space, so we are free to add elements (notice, however, that the global tangent space $TQ$ is only a vector space if $Q$ itself is a vector space!).

\section{Parallel algorithm}\label{sec:parallel-algorithm}

\todo{write about the unsuitability of sequantially getting the path, in~\cite{ferraroParallelIterativeMethod2025} after def 2.4}

In~\cite{ferraroParallelIterativeMethod2025}, an iterative algorithm that can solve the discrete Euler-Lagrange equations in proposed. The presentation here is a bit different from the original one, so that it is more in line with the extensions proposed later.

The idea of the parallel algorithm is to treat the (discrete) action as a function of which we want to find roots. I.e., given a path $\mathbf{q} = {\{q_k\}}_{k=0}^N$ and a discrete Lagrangian $\Ld$, we define its (global) residual $R$ as
\begin{equation}
	R(\mathbf{q}) = \sum_{k=1}^{N-1}
	\mdif{2} L_d (q_{k-1}, q_k) + \mdif{1} L_d (q_k, q_{k+1}) = 0
	\label{eq:residual-global}
\end{equation}
and we want to find the path $\mathbf{q}^*$ such that $r(\mathbf{q}^*) = 0$.

At a high level, two schemes are used to find this root. First, we pose a block Jacobi update step that can be done in parallel. However, the Jacobi update is still not directly computable, so we use one or more Newton-Raphson steps to compute the Jacobi update.

\subsection{Algorithm description}

The Jacobi update step consists of finding a new path $\overline{\mathbf{q}} = \{{ \overline{q}_k \}}_{k=0}^N$ that makes each summand of the residual $0$ with respect to the previous path. That is, we define the \newterm{local residual} as
\begin{equation}
	r_k(\mathbf{q}, \overline{q}) =
	\mdif{2} L_d (q_{k-1}, \overline{q}) + \mdif{1} L_d (\overline{q}, q_{k+1}) = 0
	\label{eq:residual-local}
\end{equation}
which is a term of the sum, replacing the central $q_k$ by some new value $\overline{q}$. If we find a new path $\overline{\mathbf{q}}$ such that, for all inner points $k$, $r_k(\mathbf{q}, \overline{q_k}) = 0$, that does not mean that the global residual will be $0$. However, under certain conditions, repeating this step will get the global residual arbitrarily $0$.

\todo{Cite stuff about Jacobi in~\cite[sec 4.1]{ferraroParallelIterativeMethod2025}}

Since the Lagrangian can be arbitrarily complex, this step does not have a general way to be computed. Specific kinds of Lagrangians, such as quadratic or cubic ones, have clear closed-form solutions. In lieu of knowing if the Lagrangian has one of this forms ahead of time, we use the Newton-Raphson method to solve a linearization of the problem. So, instead of consider the residual $r_k(\mathbf{q}, \overline{q})$, we consider the linearization $\rho_k$ around $q_k$:
\[
	\rho_k (\mathbf{q}, \overline{q}) = r_k(\mathbf{q}, q_k) + \mdif{2} r_k(\mathbf{q}, q_k) (\overline{q} - q_k)
\]

Notice this is the step where we have to apply the hypothesis that $Q$ is a vector space, since we are doing a linear combination based on the increment.

Finding this root is easy, since this is can be solved as the linear system:
\begin{align*}
	r_k(\mathbf{q}, q_k) + \mdif{2} r_k(\mathbf{q}, q_k) (\overline{q} - q_k) & = 0                                                                                    \\
	\overline{q} - q_k                                                        & = \left(\mdif{2} r_k(\mathbf{q}, q_k)\right)^{-1} {\left(-r_k(\mathbf{q}, q_k)\right)} \\
\end{align*}
from which we get the general Jacobi-Newton update rule:
\begin{equation}
	\overline{q} = q_k - \left(\mdif{2} r_k(\mathbf{q}, q_k)\right)^{-1} {r_k(\mathbf{q}, q_k)}
	\tag{Jacobi-Newton}
	\label{eq:jacobi-newton}
\end{equation}

In a standard discrete Lagrangian system, the terms are:

\begin{align*}
	\mdif{2} r_k (\mathbf{q}, q_k) & = \mdif{22} \Ld(q_{k-1}, q_k) + \mdif{11} \Ld(q_k, q_{k+1}) \\
	r_k (\mathbf{q}, q_k)          & = \mdif{2} \Ld(q_{k-1}, q_k) + \mdif{1} \Ld(q_k, q_{k+1})   \\
\end{align*}
(note that throughout~\cite{ferraroParallelIterativeMethod2025}, they call $\mathcal{D}_k$ what we call $\mdif{2} r_k (\mathbf{q}, q_k)$).

Equation~\eqref{eq:jacobi-newton} is finally easily computable, using any linear solver.

Of course, the question now is to verify under what conditions this algorithm produces the correct solution.

\subsection{Convergence criteria}

\todo{Cite in general}

Given a problem in the form $\nabla F = 0$, the Jacobi-Newton method converges locally if the hessian of $F$ is positive definite \tc{}. Hamilton's principle is such a formulation: extremizing the action, which quite literally means that its derivative has to be $0$. That is, $R = \nabla S$. We will call the Hessian in question (i.e., of the action) the \newterm{global Hessian matrix}, $H$.

So, given the action defined as in \autoref{eq:action-discrete} for a path $\mathbf{q}$ of length $N$
\[ 	S_d (\mathbf{q}) : Q^N \to \RR = \sum_{k=0}^{N-2} \Ld(q_k, q_{k+1})
\]
the derivative with respect to the $i$th component is, as usual,
\begin{align*}
	\mdif{i} S_d (\mathbf{q}) & = \mdif{q_i} \sum_{k=0}^{N-2} \Ld(q_k, q_{k+1})           \\
	                          & = \mdif{1} \Ld(q_i, q_{i+1}) + \mdif{2} \Ld(q_{i-1}, q_i)
\end{align*}
and then the entries of the Hessian are
\begin{align*}
	(H_{ij} (\mathbf{q})) = \mdif{q_i q_j} S_d (\mathbf{q}) & = \mdif{q_j} \left(
	\mdif{1} \Ld(q_i, q_{i+1}) + \mdif{2} \Ld(q_{i-1}, q_i)
	\right)                                                                                                                                              \\
	                                                        & = \delta_{i, j} \left( \mdif{11} \Ld(q_i, q_{i+1}) + \mdif{22} \Ld(q_{i - 1}, q_i) \right) \\
	                                                        & + \delta_{i+1,j} \mdif{12} \Ld(q_i, q_{i+1})                                               \\
	                                                        & + \delta_{i-1,j} \mdif{21} \Ld(q_{i-1}, q_i)
\end{align*}

This results in a tridiagonal matrix that can be written in the following form:
\begin{equation}
	H(\mathbf{q}) = \begin{pmatrix}
		\mathcal{D}_1       & \mathcal{C}_{1} &                         &                         &                     \\
		\mathcal{C}_{1}^{T} & \mathcal{D}_2   & \mathcal{C}_{2}         &                         &                     \\
		                    & \ddots          & \ddots                  & \ddots                  &                     \\
		                    &                 & \mathcal{C}_{N - 3}^{T} & \mathcal{D}_{N - 2}     & \mathcal{C}_{N - 2} \\
		                    &                 &                         & \mathcal{C}_{N - 2}^{T} & \mathcal{D}_{N - 1}
	\end{pmatrix}.
	\label{eq:hessian-matrix}
\end{equation}
where
\begin{equation}
	\begin{aligned}
		\mathcal{D}_k & = \mdif{22} \Ld(q_{k-1}, q_k) + \mdif{11} \Ld(q_k, q_{k+1}) \\
		\mathcal{C}_k & = \mdif{12} \Ld(q_k, q_{k+1})
	\end{aligned}
	\label{eq:hessian-terms-1}
\end{equation}
and, using local coordinates, $\mathcal{C}_k^T = \mdif{21} \Ld(q_k, q_{k+1})$ \todo{Does this actually always happen for non-smooth functions?}.

The following theorem gives one of the fundamental criteria for convergence of the parallel algorithm (\cite[Prop. 4.3, 4.5, remark 4.6]{ferraroParallelIterativeMethod2025})
\begin{theorem}[Jacobi-Newton convergence criterion]\label{thm:hessian-pd-convergence}
	If the Hessian $H(\mathbf{q})$ is positive definite, there is a neighborhood $\mathcal{D}$ around a solution $\mathbf{q}^*$ such that any path $\mathbf{q} \in \mathcal{D}$ converges to $\mathbf{q}^*$;
\end{theorem}

This condition might theoretically be good enough, although in practice it is hard to compute, since it involves checking for positive-definiteness of an $N \times N$ matrix. The most common way to do this involves the Cholesky factorization, which has an $\mathcal{O}(N^3)$ runtime.

Therefore, an easier-to-check condition is be given in terms of a \newterm{local Hessian matrix} $H_k$ (\tc{}):
\begin{equation}
	H_k(q_k, q_{k+1}) = \begin{pmatrix}
		\mathcal{A}_k   & \mathcal{C}_k \\
		\mathcal{C}_k^T & \mathcal{B}_k
	\end{pmatrix}
\end{equation}
where
\begin{equation}
	\begin{aligned}
		\mathcal{A}_k & = \mdif{11} \Ld(q_k, q_{k+1}) \\
		\mathcal{B}_k & = \mdif{22} \Ld(q_k, q_{k+1}) \\
	\end{aligned}
	\label{eq:hessian-terms-2}
\end{equation}
(note that $\mathcal{D}_k = \mathcal{B}_{k-1} + \mathcal{A}_k$).

\begin{theorem}[Local Hessian convergence~{\cite[thm. 4.1]{ferraroParallelComputingVariational}}]\label{thm:local-hessian-pd-convergence}
	Given the local Hessian matrices $H_k$, if either \begin{itemize}
		\item for all $k$, the local Hessian $H_k$ is positive definite, or
		\item \todo{I'm not sure about the second thingy}
	\end{itemize}
	then the global Hessian matrix $H$ is positive definite and the method converges locally \tc.
\end{theorem}
\begin{proof}
	Given in \tc{}
\end{proof}


This makes the checking $\mathcal{O}(N)$ in terms of the length \todo{I guess N is not necessarily ``length''} of the path. However, this condition is not necessary, and in fact in most cases of positive definite global Hessians, the conditions will not pass.

This is where a third convergence result comes into play, which gives a necessary and sufficient condition for the positive definiteness of the global Hessian matrix:

\begin{theorem}[Hessian positive definiteness criterion ({\cite[Theorem 4.12]{ferraroParallelIterativeMethod2025}})]\label{thm:hessian-pd-criterion}
	The Hessian matrix $H(\mathbf{q})$ is positive definite if and only if the matrices defined iteratively by $\Lambda_1 = \mathcal{D}_1$ and
	\[ \Lambda_i = \mathcal{D}_i - \mathcal{C}_{i-1}^T \Lambda_{i-1}^{-1} \mathcal{C}_{i-1}, \quad 2 \leq i \leq N - 1 \]
	are all positive definite.
\end{theorem}



\todo{Write it out.}
